% !TEX root = ../../main.tex

\section{Ejecución Secuencial}

El punto de comparación más simple es ejecutar el bloque en el orden monádico original, sin aprovechar \texttt{ApplicativeDo}. En ese escenario, cada statement debe esperar la finalización del anterior, aun cuando no exista una dependencia directa entre ellos.

Con el modelo de costos usado por \texttt{ApplicativeDo}, cada statement elemental tiene costo unitario. Por lo tanto, el plan secuencial del ejemplo es:

\begin{verbatim}
s1 ; s2 ; s3 ; s4
\end{verbatim}

Su costo total es la suma de los cuatro statements:

\[
1 + 1 + 1 + 1 = 4
\]

Este valor fija el baseline del capítulo. Las transformaciones posteriores se comparan contra este costo para mostrar primero la mejora obtenida por \texttt{ApplicativeDo} y luego la mejora adicional obtenida al permitir reordenamiento conmutativo.
